\documentclass[12pt]{article}
\usepackage[utf8]{inputenc}
\usepackage[spanish]{babel}
\usepackage{amsmath}
\usepackage{amssymb}
\usepackage{geometry}
\usepackage{booktabs}
\geometry{a4paper, margin=1in}

\title{Actividad 6: Gradiente Descendente en una Red Neuronal Artificial}
\author{Emilio Izquierdo Montero}
\date{\today}

\begin{document}

\maketitle

\section{Parámetros Iniciales del Problema}
Basado en el ejercicio numérico del video, definimos las siguientes condiciones y variables del sistema:
\begin{itemize}
    \item \textbf{Entradas:} $x_1 = 2$, $x_2 = 1$
    \item \textbf{Pesos iniciales (Época 1):} $w_1 = -0.5$, $w_2 = 0.5$
    \item \textbf{Sesgo (Bias):} $b = 0$
    \item \textbf{Valor esperado (Target):} $y = 0.6$
    \item \textbf{Tasa de aprendizaje (Learning rate):} $\eta = 0.4$
    \item \textbf{Función de activación:} Sigmoide $\sigma(h) = \frac{1}{1 + e^{-h}}$
\end{itemize}

\section{Fórmulas Utilizadas}
\begin{enumerate}
    \item \textbf{Suma neta:} $h = w_1x_1 + w_2x_2 + b$
    \item \textbf{Salida predicha:} $\hat{y} = \sigma(h) = \frac{1}{1 + e^{-h}}$
    \item \textbf{Error residual simple:} $e_{residual} = y - \hat{y}$
    \item \textbf{Error cuadrático (Loss):} $E = (y - \hat{y})^2$
    \item \textbf{Término de error del gradiente:} $\delta = (y - \hat{y}) \cdot \hat{y} \cdot (1 - \hat{y})$
    \item \textbf{Incremento de pesos:} $\Delta w_i = \eta \cdot \delta \cdot x_i$
    \item \textbf{Actualización de pesos:} $w_i^{nuevo} = w_i + \Delta w_i$
\end{enumerate}

\newpage

\section{Desarrollo de las Épocas}

\subsection{Época 1}
\begin{itemize}
    \item \textbf{Forward Pass:}
    \[ h = (-0.5 \times 2) + (0.5 \times 1) = -1.0 + 0.5 = -0.5 \]
    \[ \hat{y} = \sigma(-0.5) = 0.377 \]
    \item \textbf{Cálculo de Errores:}
    \[ e_{residual} = 0.6 - 0.377 = 0.223 \]
    \[ E = (0.223)^2 = 0.049 \]
    \item \textbf{Gradiente y Actualización:}
    \[ \delta = 0.223 \times 0.377 \times (1 - 0.377) = 0.052 \]
    \[ \Delta w_1 = 0.4 \times 0.052 \times 2 = 0.041 \implies w_1^{nuevo} = -0.5 + 0.041 = -0.459 \]
    \[ \Delta w_2 = 0.4 \times 0.052 \times 1 = 0.020 \implies w_2^{nuevo} = 0.5 + 0.020 = 0.520 \]
\end{itemize}

\subsection{Época 2}
\begin{itemize}
    \item \textbf{Forward Pass:}
    \[ h = (-0.459 \times 2) + (0.520 \times 1) = -0.918 + 0.520 = -0.398 \]
    \[ \hat{y} = \sigma(-0.398) = 0.401 \]
    \item \textbf{Cálculo de Errores:}
    \[ e_{residual} = 0.6 - 0.401 = 0.199 \]
    \[ E = (0.199)^2 = 0.0396 \]
    \item \textbf{Gradiente y Actualización:}
    \[ \delta = 0.199 \times 0.401 \times (1 - 0.401) = 0.0478 \]
    \[ \Delta w_1 = 0.4 \times 0.0478 \times 2 = 0.0382 \implies w_1^{nuevo} = -0.459 + 0.0382 = -0.421 \]
    \[ \Delta w_2 = 0.4 \times 0.0478 \times 1 = 0.0191 \implies w_2^{nuevo} = 0.520 + 0.0191 = 0.539 \]
\end{itemize}

\subsection{Época 3}
\begin{itemize}
    \item \textbf{Forward Pass:}
    \[ h = (-0.421 \times 2) + (0.539 \times 1) = -0.842 + 0.539 = -0.303 \]
    \[ \hat{y} = \sigma(-0.303) = 0.425 \]
    \item \textbf{Cálculo de Errores:}
    \[ e_{residual} = 0.6 - 0.425 = 0.175 \]
    \[ E = (0.175)^2 = 0.0306 \]
    \item \textbf{Gradiente y Actualización:}
    \[ \delta = 0.175 \times 0.425 \times (1 - 0.425) = 0.0428 \]
    \[ \Delta w_1 = 0.4 \times 0.0428 \times 2 = 0.0342 \implies w_1^{nuevo} = -0.421 + 0.0342 = -0.387 \]
    \[ \Delta w_2 = 0.4 \times 0.0428 \times 1 = 0.0171 \implies w_2^{nuevo} = 0.539 + 0.0171 = 0.556 \]
\end{itemize}

\newpage

\section{Respuestas del Cuestionario}

\subsection{¿Cuál fue el Error Residual en la Época 1, 2 y 3?}
El comportamiento del error residual simple disminuyó consistentemente y condisderablemente en cada iteración:
\begin{itemize}
    \item \textbf{Época 1:} $0.223$
    \item \textbf{Época 2:} $0.199$
    \item \textbf{Época 3:} $0.175$
\end{itemize}

\subsection{¿Cómo se comportó la salida predicha ($\hat{y}$) en cada época? ¿La red está aprendiendo?}
La salida predicha evolucionó de la siguiente manera:
\begin{itemize}
    \item \textbf{Época 1:} $\hat{y} = 0.377$
    \item \textbf{Época 2:} $\hat{y} = 0.401$
    \item \textbf{Época 3:} $\hat{y} = 0.425$
\end{itemize}
\textbf{Conclusión:} \textbf{Sí, está aprendiendo.} Dado que nuestro target es $y = 0.6$, el incremento progresivo de $\hat{y}$ hacia ese valor demuestra que el algoritmo de optimización está minimizando la función de costo 
\subsection{¿Cómo se comportaron los pesos $w_i$ en cada época?}
Los pesos se reajustaron dinámicamente según la magnitud de sus entradas :
\begin{itemize}
    \item El peso \textbf{$w_1$} (asociado a $x_1 = 2$) aumentó su valor  desde $-0.5$ hasta $-0.387$, aproximándose hacia el cero.
    \item El peso \textbf{$w_2$} (asociado a $x_2 = 1$) se incrementó positivamente de manera constante, pasando de $0.5$ a $0.556$.
\end{itemize}
Ambos pesos cambiaron en la dirección opuesta al gradiente.

\end{document}
